\section{Русскоязычный бенчмарк RAGBench-RU}
\label{sec:dataset}

\subsection{Исходный RAGBench и отбор поддатасетов}
Оригинальный RAGBench~\cite{friel2024ragbench} включает 12 поддатасетов,
 относящихся к пяти предметным доменам. В рамках настоящей работы для построения 
 русскоязычного аналога был сформирован поднабор из 7 поддатасетов, обеспечивающий 
 доменное разнообразие и достаточный объём данных для обучения и оценки моделей:

\begin{itemize}
  \item \textbf{общие знания}: \texttt{hotpotqa}, \texttt{hagrid}, \texttt{expertqa};
  \item \textbf{биомедицина}: \texttt{covidqa};
  \item \textbf{клиентская поддержка}: \texttt{delucionqa};
  \item \textbf{юридические контракты}: \texttt{cuad};
  \item \textbf{финансы}: \texttt{finqa}.
\end{itemize}

Оставшиеся поддатасеты (\texttt{msmarco}, \texttt{pubmedqa}, \texttt{emanual},
 \texttt{techqa}, \texttt{tatqa}) в данную версию корпуса не включались. 
 Полученный русскоязычный корпус содержит около 31~тыс. примеров, разделённых
 на обучающую, валидационную и тестовую выборки; этого оказалось достаточно для
 основных экспериментов настоящей работы.

\subsection{Схема одного примера}

Каждая запись RAGBench/RAGBench-RU описывается JSON-подобной структурой со
следующими ключами: вопрос (\texttt{question}), список документов в виде списков 
предложений (\texttt{documents\_sentences}), ответ (\texttt{response}), 
а также GPT-4-разметка в виде списков ключей релевантных и использованных предложений 
и булева флага \texttt{adherence\_score}. Пример записи из \texttt{delucionqa} 
в сокращённом виде:

\begin{lstlisting}[basicstyle=\ttfamily\scriptsize,breaklines=true]
{
  "question": "Does the 2019 Jeep Cherokee have a backup camera?",
  "documents_sentences": [
    [["0a", "Jeep Cherokee..."], ["0b", "ParkView backup camera..."]]
  ],
  "response": "Yes, the 2019 Jeep Cherokee includes a backup camera...",
  "all_relevant_sentence_keys": ["0b"],
  "all_utilized_sentence_keys": ["0b"],
  "adherence_score": true
}
\end{lstlisting}

Для модели-оценщика из таких записей формируется последовательность 
$[q][\text{SEP}][\text{документы}][\text{SEP}][a]$ (см. раздел~\ref{sec:problem}); 
метки \texttt{relevant\_sentence\_keys} и \texttt{utilized\_sentence\_keys} 
разворачиваются в бинарные маски на токенах, \texttt{adherence\_score} - на токенах ответа.

\subsection{Пайплайн перевода}
\label{sec:translation_pipeline}

Ядро перевода - асинхронный конвейер, построенный на Qwen2.5-72B-Instruct~\cite{qwen2}.
Ключевые инженерные решения:

\begin{enumerate}
  \item \textbf{Отдельные \glslink{промпт}{промпты} для каждого поля.} Вопросы и ответы переводятся
  одним запросом на поле; \texttt{documents\_sentences} - построчно, с
  фиксированным порядком предложений, чтобы сохранить корректность ключей
  \texttt{a/b/c/\dots}, к которым привязана разметка.
  \item \textbf{Инструкция-каркас для документов.} Промпт явно запрещает: (а)~менять 
  число предложений в документе, (б)~добавлять нумерацию или Markdown, (в)~переводить
   английские собственные имена, числа и URL. Добавление таких правил снизило долю формальных
    рассогласований на этапе склейки примерно на порядок.
  \item \textbf{Пакетная отправка.} Отправка идёт батчами по 16-32 примера на 
  один GPU Qwen через \texttt{asyncio} и семафор по числу параллельных запросов; 
  это типовой приём, использующийся как в оригинальной статье RAGBench, так и 
  в ARES~\cite{saad2023ares}.
  \item \textbf{Двухстадийная постобработка.} После получения ответа модели проверяются:
   (а)~равенство числа предложений до и после перевода, (б)~сохранность ключей предложений
    в \texttt{all\_relevant\_sentence\_keys} / \texttt{all\_utilized\_sentence\_keys}. 
    Примеры, не прошедшие проверку, переотправляются с более строгим системным промптом;
    при повторной неудаче отбрасываются (итоговая доля отброшенных $<$ 1\%).
\end{enumerate}

Формулировки промптов, использованных на финальной версии, приведены в приложении.

\subsection{LLM-разметка и контроль качества}
\label{sec:llm_eval}

Проверка качества перевода проведена в три этапа.

\textbf{1. Автоматические счётчики.} Для каждого примера сопоставлены: число предложений 
в \texttt{documents\_sentences} до/после, длина (в символах и токенах) вопроса/ответа, 
наличие всех исходных ключей разметки. Проверку прошли 99\% примеров; неуспешные отправлены
 на повторный перевод.

\textbf{2. Выборочный визуальный отсмотр.} Автор работы просматривал случайные примеры
 из каждого поддатасета на предмет очевидных сбоев перевода: обрывы предложений, смешение
  русского и английского, потеря или дублирование предложений, нарушение ключей
   \texttt{a/b/c/\dots}. Результаты отсмотра использовались итеративно - для уточнения
    системного промпта и инструкции-каркаса перевода (раздел~\ref{sec:translation_pipeline});
     формальной разметки ошибок и количественных оценок адекватности перевода на этом этапе
      не велось.

\textbf{3. Консистентность разметки.} Так как русскоязычные примеры используют те же
 ключи предложений, что и английские оригиналы, разметка \texttt{all\_relevant\_sentence\_keys}
 и \texttt{all\_utilized\_sentence\_keys} переносится без изменений. Этот приём известен по
 работам, использующим параллельные корпуса: машинный перевод не меняет структуру, только
 содержимое предложений.

Следует явно отметить: ручная проверка перевода носителями языка
 (независимый перевод и подсчёт согласия) в данной работе не проводилась. Это ограничение
  отмечено в разделе~\ref{sec:conclusion} и вынесено в направления дальнейшей работы.

\subsection{Статистика корпуса}

Итоговые размеры сплитов по 7 поддатасетам представлены в 
таблице~\ref{tab:dataset_stats}. Распределение по доменам сбалансировано:
 общие знания (\texttt{hotpotqa}, \texttt{hagrid}, \texttt{expertqa}) составляют 
 около 60\% корпуса, остальное приходится на специализированные поддатасеты.

\begin{table}[H]
\centering
\small
\caption{Размеры сплитов RAGBench-RU по поддатасетам}
\label{tab:dataset_stats}
\begin{tabular}{lcccc}
\toprule
поддатасет & train & val & test & домен \\
\midrule
hotpotqa   & 1880 & 424 & 390 & общие знания \\
hagrid     &  2890 &  322 &  1322 & общие знания \\
expertqa   &  1490 &  196 &  188 & общие знания \\
covidqa    &  1250 &  267 &  246 & биомедицина \\
delucionqa &   1460 &  182 &  184 & клиентская поддержка \\
cuad       &  749 &  269 &  305 & юриспруденция \\
finqa      &  12500 &  1770 &  2290 & финансы \\
\midrule
\textbf{итого} & \textbf{22 219} & \textbf{3 430} & \textbf{4 925} & - \\
\bottomrule
\end{tabular}
\end{table}

Распределения длин (в токенах \texttt{deepvk/RuModernBERT-base}) 
существенно различаются по поддатасетам: медиана общей длины входа (вопрос+контекст+ответ) 
варьируется от $\sim$250 токенов (\texttt{hotpotqa}, \texttt{hagrid}) 
до $\sim$1200--1500 (\texttt{cuad}, \texttt{finqa}), с хвостом до 8~тыс. токенов 
в единичных примерах \texttt{cuad}. Это мотивирует эксперимент по длине контекста
 (раздел~\ref{sec:experiments}): одни поддатасеты фактически требуют длины 4096-8192,
  другие полностью укладываются в 512.

\subsection{Публикация}

Итоговый корпус опубликован на платформе Hugging Face Datasets под именем
\texttt{CMCenjoyer/ragbench-ru}
\footnote{\href{https://huggingface.co/datasets/CMCenjoyer/ragbench-ru}{https://huggingface.co/datasets/CMCenjoyer/ragbench-ru}}.
Публикация включает все 7 поддатасетов с разбиением на обучающую, 
валидационную и тестовую выборки. Лицензия корпуса совпадает с лицензией
  исходного RAGBench. В файле \texttt{README} указаны источник данных и 
  способ перевода (Qwen2.5-72B, машинный перевод с помощью промпта без ручной постредакции), 
  что позволяет учитывать возможный шум перевода при интерпретации результатов.
